\section{Conclusion}
\label{sec:conclusion}

This study asks how the definition of demonstration difficulty affects
curriculum behavior cloning for robotic manipulation. We compare Task, Expert,
and out-of-fold Learner difficulty under five sampling strategies while
matching the policy, initialization protocol, optimizer, update budget,
dataset, and evaluation tasks. The metrics induce substantially different
rankings, and the three easy-to-hard schedules do not consistently outperform
Uniform.

Reverse Learner attains the highest observed mean pick-and-place success, but
the result is seed-sensitive and is not accompanied by faster development AUC.
The near-floor stacking study provides no decisive cross-task replication.
Because staged eligibility couples ordering and exposure, the experiments
support a cautious conclusion: difficulty definitions and sampling schedules
must be evaluated with matched baselines and closed-loop outcomes rather than
ranking intuition or offline loss alone.
